\section{Identification of the fixed boundary descent}
\label{bm-section}

We discharge the explicit modular-membership premise for the boundary
curve of Section~\ref{bs-section}. The argument uses one fixed
Tate-module prime and a finite character comparison, followed by
the compatibility of an already constructed modular form. It does not
assume that unrelated choices of descent at different primes are
compatible. The resulting orbit identification combines ordinary
modularity theorems with the independently repeated exact newspace
enumeration; it is not a Lean formalization of these inputs.

Write $K=\mathbb Q(r)$, $r^2=-3$, and
\[
 E_0:\quad Y^2=X^3+12X^2+6(3+r)X.
\]
Theorem~\ref{bs-noncm} proves geometric non-CM. The same explicit
two-isogeny as in the actual family identifies its quotient with the
$(-2)$-twist of its conjugate, so $E_0$ is a Q-curve. Fix the
order-four character $\chi$ used in the earlier descent. Its square
is $\psi_3|_{G_K}$, and its conjugation relation cancels the character
of that isogeny.

\begin{lemma}[Finite-character comparison]\label{bm-character}
Let $H$ be a finite-index normal subgroup of $G$. Let $\rho_1,\rho_2$
be representations over an algebraically closed characteristic-zero
field whose restrictions to $H$ are isomorphic and absolutely
irreducible. Then $\rho_2\simeq\rho_1\otimes\nu$ for a character
$\nu$ of the finite group $G/H$.
\end{lemma}
\begin{proof}
Identify the two restrictions. For $g\in G$, the matrix
$\rho_2(g)\rho_1(g)^{-1}$ commutes with their common action of $H$,
by normality. Schur's lemma makes it scalar, say $\nu(g)$. The two
representation laws give $\nu(gg')=\nu(g)\nu(g')$, and $\nu|_H=1$.
For continuous Galois representations this is a finite continuous
character; over $G_{\mathbb Q}$ it is a Dirichlet character.
\end{proof}

\begin{theorem}[Modularity of the specified boundary descent]
\label{bm-modular}
Either extension of $V_{17}(E_0)\otimes\chi$ to $G_{\mathbb Q}$
is the representation of a weight-two newform with nebentypus
$\psi_3$.
\end{theorem}
\begin{proof}
The conjugation relation and Schur normalization in the proof of
\cite[Theorem~4.2]{PacettiVillagra2022QCurves} give the extension.
Its two choices differ by $\psi_{-3}$. Here only the general
representation argument is used, not that source's equation-specific
conductor formula.

Ribet's Theorem~6.1 and its construction give a primitive abelian
variety $A$ of GL$_2$-type over $\mathbb Q$ and a finite Galois
extension $L/\mathbb Q$ such that $A_L$ is isogenous to a power of
$E_0$; see the
\href{https://math.berkeley.edu/~ribet/Articles/korea.pdf}{author's
\emph{Abelian varieties over $\mathbb Q$ and modular forms}},
Section~6. Enlarge $L$ to contain $K$ and kill $\chi$.
The completed modularity theorem for GL$_2$-type abelian varieties,
\cite[Corollary~10.2(i)]{KhareWintenberger2009Serre}, makes $A$
modular. Choose a two-dimensional coefficient-field component
$\rho_A$ of its Tate module at $17$. It comes from a weight-two
newform $f_A$.

On $G_L$, this component is isomorphic to $V_{17}(E_0)$ after
scalar extension: the whole Tate module is a sum of copies and the
chosen component has dimension two. Faltings' semisimplicity and
endomorphism theorem, together with
$\operatorname{End}_L(E_0)\otimes\mathbb Q=\mathbb Q$, make
this restriction absolutely irreducible. The specified extension
$\widetilde\rho$ has the same restriction, since $\chi$ is killed
there. Lemma~\ref{bm-character} gives
\[
 \widetilde\rho\simeq\rho_A\otimes\nu
\]
for a finite Dirichlet character $\nu$. The primitive newform
associated to $f_A\otimes\nu$ still has weight two, and realizes
$\widetilde\rho$. In particular it is odd: a finite character
has square $1$ on complex conjugation. On $G_K$ the determinant
is $\psi_3$ times the cyclotomic character; oddness gives equality
on complex conjugation and therefore on $G_{\mathbb Q}$. Its
nebentypus is exactly $\psi_3$.
\end{proof}

The multiplicative-prime large-image result for positive seeds has
not been used; it would not apply to $F(0,1)=1$. Nor has any
nonzero Tate twist been introduced.

\begin{lemma}[Compatibility supplied by the identified modular form]
\label{bm-compatible}
For the form $f$ in Theorem~\ref{bm-modular} and every auxiliary
prime $\ell$, after fixing compatible coefficient embeddings,
\[
 \rho_{f,\ell}|_{G_K}
       \simeq V_\ell(E_0)\otimes\chi_\ell.
\]
\end{lemma}
\begin{proof}
At every prime of $K$ outside a finite set, the restricted
Frobenius polynomials agree by the isomorphism at $17$. They are
polynomials in algebraic numbers, and the embedding into
$\overline{\mathbb Q}_{17}$ is injective. Compatibility for the
elliptic curve, finite character and already constructed modular
form gives equality at any other auxiliary prime. Chebotarev and
semisimplicity then give the displayed isomorphism.
\end{proof}

\begin{proposition}[The actual characteristic-zero level]
\label{bm-level}
The newform in Theorem~\ref{bm-modular} has exact level
\[
 N=2^e3^2,\qquad 2\le e\le6.
\]
\end{proposition}
\begin{proof}
The local Tate branches hold at $(a,b)=(0,1)$, and can be checked
directly. At the inert prime over $2$, translate $X$ by $1+r$ and
$Y$ by $4$. The new coefficients are
\[
 a'_1=0,\quad a'_2=15+3r,\quad a'_3=8,\quad
 a'_4=36+36r,\quad a'_6=-48+48r.
\]
Their relevant valuations are $1,3,3,5$. For example
$a'_6=96(\zeta-1)$ with $\zeta=(1+r)/2$ and $\zeta-1$ a unit
at $2$. The discriminant valuation is $12$. The first Tate cubic
has the form $Z^2(Z+b)$ with $b$ a unit; the first starred
quadratic has double root zero, and the next has unit linear
coefficient $a'_4/8$. Thus the minimal type is $I_2^*$ and its
conductor exponent is $6$, by the earlier reviewed Tate algorithm.

At $r$ over $3$, with $v_r(3)=2$, the original coefficients have
$v_r(a_2)=2$, $v_r(a_4)=3$, $a_1=a_3=a_6=0$ and
$v_r(\Delta)=9$. The successive multiple roots are zero and
$a_4/r^3$ is a unit. The minimal type is $III^*$, of conductor
exponent $2$. At all primes above $q>3$ the original curve has good
reduction.

The character $\chi$ is unramified at $3$ and has conductor
exponent $3$ at $2$. At $2$ the elliptic representation has finite
inertia and unramified determinant. Every inertia-subgroup fixed
space has dimension zero or two: a finite determinant-one matrix
fixing a nonzero vector must be the identity. The conductor $6$
therefore consists of tame codimension $2$ and Swan conductor $4$,
so its final upper break is $2$. The character also has upper break
$2$. Their tensor product is trivial on all upper groups with
index greater than $2$, giving conductor at most $6$.
Unramified restriction from $\mathbb Q_2$ to $K_2$ preserves the
exponent. The determinant $\psi_3$ has conductor exponent $2$
there, giving the lower bound $2$.

At $3$ the restricted representation is tame with no inertia
invariants. The extension $K_3/\mathbb Q_3$ is tame, so its wild
inertia groups agree. The descent is tame, and its inertia
invariants are contained in those of the restriction. Its conductor
exponent is consequently exactly $2$ for either descent.

For $q>3$ choose $\ell\ne q$ and apply
Lemma~\ref{bm-compatible}. The elliptic curve and character are
unramified at $q$, and $K/\mathbb Q$ is unramified there, so the
descent is unramified. In particular the rational prime $17$ can
be checked using $\ell=19$. Equality of characteristic-zero
newform conductor and the conductor of its representation away
from the auxiliary prime now gives the asserted level.
\end{proof}

\begin{theorem}[The non-CM level-$576$ orbit]\label{bm-orbit}
Together with the complete exact newspace enumeration recorded in
Section~\ref{bs-section}, the preceding propositions identify either
fixed boundary descent with the unique non-CM Galois orbit of weight
two, nebentypus $\psi_3$, and level $576$.
\end{theorem}
\begin{proof}
The form cannot have CM. Over a finite extension a CM form's
two-dimensional representation decomposes into characters, whereas
the Tate representation of $E_0$ remains absolutely irreducible.
Lemma~\ref{bm-compatible} gives a contradiction after enlarging
the finite extension to kill $\chi$.

Among the five levels of Proposition~\ref{bm-level}, the complete
enumeration has exactly two non-CM orbits, at $288$ and $576$.
Their rational coefficients $a_{13}$ are $4$ and $-4$ respectively.
Lemma~\ref{bs-character-sign} gives $a_{13}=-4$. Since thirteen
splits in $K$, the ambiguity $\psi_{-3}$ between the two descents
does not change that trace. The level-$576$ orbit is therefore the
unique possibility.
\end{proof}

This uses a proved membership theorem followed by a complete finite
enumeration. It is not a conclusion of equality of representations
from a few matching coefficients. The finite computation and its
independent exact replay remain a separate, explicit dependency.

For the specified descent, this discharges the premise in
Proposition~\ref{bs-conditional-mazur}: those particular finite-state
products have a zero factor at every auxiliary prime. It supplies
no global power equation for their positive local shadows. Additional
global covers or other restrictions can still eliminate actual
integer seeds. In particular the independent effective affine-slice
argument limits the consecutive shadow family globally.

The complete ordinary comparison, local-conductor calculation and
orbit-identification argument passed two independent reviews. Both
reviewers opened Ribet's construction and Khare--Wintenberger's
GL$_2$-type modularity corollary. No new Lean theorem is asserted.
